% !TeX root = ../main.tex
\chapter{Mô tả dataset}
\label{chap:dataset}

\section{Nguồn dữ liệu}

Nhóm sử dụng Breast Cancer Wisconsin (Diagnostic), do Wolberg, Mangasarian,
Street và Street công bố trên UCI Machine Learning Repository
\cite{uci-breast-cancer}. Notebook tải bản dữ liệu được đóng gói trong
scikit-learn bằng \texttt{load\_breast\_cancer(as\_frame=True)}; vì vậy không
có đường dẫn tuyệt đối hoặc file dữ liệu riêng phụ thuộc máy của thành viên.

Các feature được tính từ ảnh số hóa của phép chọc hút kim nhỏ một khối u vú.
Chúng mô tả đặc trưng của nhân tế bào như bán kính, texture, perimeter, area,
smoothness, compactness, concavity, concave points, symmetry và fractal
dimension. Mỗi nhóm đo có giá trị mean, standard error và worst, tổng cộng 30
feature đầu vào.

\begin{table}[H]
\centering
\caption{Tóm tắt Breast Cancer Wisconsin (Diagnostic)}
\label{tab:dataset-summary}
\begin{tabular}{p{5.0cm}p{8.5cm}}
\toprule
Thuộc tính & Giá trị \\
\midrule
Nguồn & UCI Machine Learning Repository / scikit-learn \\
Số mẫu & 569 \\
Số feature & 30 feature số thực \\
Biến mục tiêu & Chẩn đoán khối u \\
Nhãn lớp & malignant (0), benign (1) \\
Phân bố lớp & 212 malignant (37,26\%); 357 benign (62,74\%) \\
Giá trị thiếu & 0 \\
Dòng feature trùng hoàn toàn & 0 \\
\bottomrule
\end{tabular}
\end{table}

\section{Tiền xử lý}

Notebook kiểm tra missing value và duplicate trước khi chia dữ liệu. Không cần
imputation vì tập dữ liệu không có giá trị thiếu; không cần encoding vì toàn bộ
feature là số. Nhóm giữ nguyên 30 feature để các threshold của cây vẫn gắn với
đặc trưng gốc. Decision Tree không phụ thuộc vào khoảng cách Euclid nên không
cần standardization.

Dataset có chênh lệch lớp ở mức vừa phải. Nhóm dùng stratified split để duy trì
tỷ lệ lớp giữa training và testing, đồng thời báo cáo recall riêng cho lớp
malignant để accuracy tổng thể không che giấu số ca ác tính bị bỏ sót.

\section{Lý do dataset phù hợp}

Đây là bài toán phân loại nhị phân có nhãn rõ ràng và các feature số hỗ trợ
phép chia theo ngưỡng. Quy mô 569 mẫu đủ để minh họa quá trình xây dựng cây và
cross-validation nhưng vẫn cho phép trực quan hóa. Các rule như
\texttt{worst radius $\leq$ 16.80} cũng có thể đọc và liên hệ với đặc trưng
hình thái, phù hợp với yêu cầu phân tích cây của Lab 2.
